\section{Pipeline and Rendering Implementation}
\label{sec:pipeline}
\subsection{Checked Pipeline Interface}
% Role: source-backed overview; detailed upstream methods await owner review.
The repository connects six CUDA stages: preprocessing (S1), keypoint detection
(S2), description (S3), matching (S4), homography estimation (S5), and rendering
(S6). S1--S3 process both images before S4 matches their descriptors. S5 returns
a homography $H$ mapping B to A in original-image coordinates. Upstream scale
metadata converts an estimate from resized coordinates before rendering.
Detailed descriptions of S1, S2/S3, and S4/S5 require their respective owners'
review. The following subsections describe the S6 code checked for this draft.

\subsection{S6 Geometry, Sampling, and Validity}
\label{sec:s6-geometry}
% Role: motivation and representation; validity is independent of RGB intensity.
S6 takes original RGB8 images and a row-major, double-precision homography
$H_{B\rightarrow A}$. The caller owns pitched image views, workspace, and a CUDA
stream. Geometry preparation rejects non-finite or singular transforms,
a projective horizon crossing B, and a canvas exceeding its configured budget.
It normalizes $H$ by its largest-magnitude coefficient and does not require
$h_{33}=1$.

% Role: forward geometry; match integer pixel-centre and origin conventions.
Pixel centres have integer coordinates. The canvas encloses A and the projected
corner centres of B using snapped floor/ceiling bounds with an inclusive $+1$
extent. The origin $(x_0,y_0)$ is expressed in A coordinates. For canvas pixel
$(x,y)$, sampling uses
\begin{equation}
 p_A=(x+x_0,y+y_0,1)^\mathsf{T},\qquad p_B\sim H^{-1}p_A.
\label{eq:inverse-warp}
\end{equation}
The renderer checks the projective denominator and source-coordinate validity
before bilinear sampling. Coordinates use FP64; colour/filter computations use
FP32. Neighbours at a valid last row or column are clamped. Clamping does not
make an out-of-source coordinate valid.

% Role: output interpretation and invariant.
Each source has an explicit validity mask independent of colour: black pixels
can be valid. The output mask is their union, encoded as 0 or 1. A pixel covered
by one source retains that source's warped colour. Pixels outside both sources
have zero RGB and mask. These rules also apply after multi-band reconstruction,
so filtering does not expand reported coverage.

\subsection{Feather Blending}
% Role: forward formula; source-normalized geometric weights are project policy.
For valid coordinate $(u_i,v_i)$ in source $i$ of size $W_i\times H_i$, define
\begin{equation}
 d_i=\min\!\left\{\frac{u_i+1/2}{W_i},\frac{W_i-1/2-u_i}{W_i},
                 \frac{v_i+1/2}{H_i},\frac{H_i-1/2-v_i}{H_i}\right\}.
\end{equation}
In the overlap, feather output is
$C_F=(d_A C_A+d_B C_B)/(d_A+d_B)$. The half-pixel terms keep edge weights
positive, including single-pixel dimensions. This source-normalized rule is
a project convention, not a canvas-mask distance transform or an assertion of
equivalence to a library feather implementation. The CUDA path fuses inverse
warping and feather composition and requires no temporary workspace.

\subsection{Multi-Band Blending}
% Role: ownership and extension; describe actual boundary handling.
Multi-band ownership assigns overlap pixels to the source with larger $d_i$,
breaking ties in favour of A. A single valid source owns its pixels. Before
building colour pyramids, a pull--push extension fills invalid colour samples
for filtering only. Colour numerators and validity weights are reduced to
$1\times1$; valid normalized values are retained while missing values are expanded
from coarser levels. Original validity and ownership masks remain unchanged.
Both requested multi-band depths use the same extension procedure.

% Role: pyramid construction and reconstruction.
Let $R$ denote filtering with the separable kernel $[1,4,6,4,1]/16$ followed by
even-coordinate subsampling to ceiling-half shape. Expansion $E$ inserts zeros,
applies the two-dimensional filter with gain four, and crops to the parent
shape. Borders use reflect-101, with an explicit rule for a dimension of length
one. For source $i$, Gaussian colours are $P_i^0=F_i$,
$P_i^{\ell+1}=R(P_i^\ell)$, where $F_i$ is the extended image. Gaussian ownership
weights use $U_i^{\ell+1}=R(U_i^\ell)$. Laplacian bands are
\begin{equation}
 L_i^\ell=P_i^\ell-E(P_i^{\ell+1})\quad(\ell<L),\qquad L_i^L=P_i^L.
\end{equation}
Each blended band is
\begin{equation}
 Z^\ell=\frac{U_A^\ell L_A^\ell+U_B^\ell L_B^\ell}{U_A^\ell+U_B^\ell},
\end{equation}
with zero for a zero denominator. Reconstruction starts at $C^L=Z^L$ and uses
$C^\ell=Z^\ell+E(C^{\ell+1})$. The modes request $L=3$ or $L=5$
\emph{reductions}, corresponding to four or six stored levels on a sufficiently
large canvas. Effective depth stops at $1\times1$.

% Role: numeric/resource trade-offs and unimplemented optimizations.
Only overlap pixels use reconstructed colour. Quantization is
$\lfloor\operatorname{clamp}(255C,0,255)+1/2\rfloor$; signed Laplacian bands are
not clamped prematurely. Current kernels evaluate the $5\times5$ filter directly,
rather than claiming a separable-pass optimization. S6 performs no allocation
or synchronization inside its asynchronous render call. The caller queries
workspace size, respects alignment and alias restrictions, and checks stream
completion before consuming output. Optional union/overlap counters have an
explicit device-buffer lifetime and cost. The renderer implements neither gain
correction, gamma linearization, content-aware seam optimization, nor warp caching.
